%%% The main file. It contains definitions of basic parameters and includes all other parts.

%% Settings for single-side (simplex) printing
% Margins: left 40mm, right 25mm, top and bottom 25mm
% (but beware, LaTeX adds 1in implicitly)
\documentclass[12pt,a4paper]{report}
\setlength\textwidth{145mm}
\setlength\textheight{247mm}
\setlength\oddsidemargin{15mm}
\setlength\evensidemargin{15mm}
\setlength\topmargin{0mm}
\setlength\headsep{0mm}
\setlength\headheight{0mm}
% \openright makes the following text appear on a right-hand page
\let\openright=\clearpage

%% Settings for two-sided (duplex) printing
% \documentclass[12pt,a4paper,twoside,openright]{report}
% \setlength\textwidth{145mm}
% \setlength\textheight{247mm}
% \setlength\oddsidemargin{14.2mm}
% \setlength\evensidemargin{0mm}
% \setlength\topmargin{0mm}
% \setlength\headsep{0mm}
% \setlength\headheight{0mm}
% \let\openright=\cleardoublepage

%% Generate PDF/A-2u
\usepackage[a-2u]{pdfx}

%% Character encoding: usually latin2, cp1250 or utf8:
\usepackage[utf8]{inputenc}

%% Prefer Latin Modern fonts
\usepackage{lmodern}

%% Further useful packages (included in most LaTeX distributions)
\usepackage{amsmath}        % extensions for typesetting of math
\usepackage{amsfonts}       % math fonts
\usepackage{amsthm}         % theorems, definitions, etc.
\usepackage{bbding}         % various symbols (squares, asterisks, scissors, ...)
\usepackage{bm}             % boldface symbols (\bm)
\usepackage{graphicx}       % embedding of pictures
\usepackage{fancyvrb}       % improved verbatim environment
\usepackage{natbib}         % citation style AUTHOR (YEAR), or AUTHOR [NUMBER]
\usepackage[nottoc]{tocbibind} % makes sure that bibliography and the lists
			    % of figures/tables are included in the table
			    % of contents
\usepackage{dcolumn}        % improved alignment of table columns
\usepackage{booktabs}       % improved horizontal lines in tables
\usepackage{paralist}       % improved enumerate and itemize
\usepackage{xcolor}         % typesetting in color

%%% Basic information on the thesis

% Thesis title in English (exactly as in the formal assignment)
\def\ThesisTitle{Conjugate Gradient Method for Solving Singular Systems}

% Author of the thesis
\def\ThesisAuthor{Janika Pokorná}

% Year when the thesis is submitted
\def\YearSubmitted{2024}

% Name of the department or institute, where the work was officially assigned
% (according to the Organizational Structure of MFF UK in English,
% or a full name of a department outside MFF)
\def\Department{Department of Numerical Mathematics}

% Is it a department (katedra), or an institute (ústav)?
\def\DeptType{Department}

% Thesis supervisor: name, surname and titles
\def\Supervisor{RNDr. Jan Papež, Ph.D.}

% Supervisor's department (again according to Organizational structure of MFF)
\def\SupervisorsDepartment{Department of Numerical Mathematics}

% Study programme and specialization
\def\StudyProgramme{Mathematical Modeling}
\def\StudyBranch{??}

% An optional dedication: you can thank whomever you wish (your supervisor,
% consultant, a person who lent the software, etc.)
\def\Dedication{%
Dedication.
}

% Abstract (recommended length around 80-200 words; this is not a copy of your thesis assignment!)
\def\Abstract{%
We rigorously derive the Conjugate Gradient (CG) method for regular systems. Then, using orthogonal decomposition of vectors, we show the conditions necessary for the CG method to work on singular systems and we prove in which cases the CG method will not work at all. We analyze the algorithm mathematically to disprove possible convergence for perturbed right-hand side. We support our results with numerical experiments.
}

% 3 to 5 keywords (recommended), each enclosed in curly braces
\def\Keywords{%
{Conjugate} {Gradients} {Singular systems} {Orthogonal decomposition}
}

%% The hyperref package for clickable links in PDF and also for storing
%% metadata to PDF (including the table of contents).
%% Most settings are pre-set by the pdfx package.
\hypersetup{unicode}
\hypersetup{breaklinks=true}

% Definitions of macros (see description inside)
%%% This file contains definitions of various useful macros and environments %%%
%%% Please add more macros here instead of cluttering other files with them. %%%

%%% Minor tweaks of style

% These macros employ a little dirty trick to convince LaTeX to typeset
% chapter headings sanely, without lots of empty space above them.
% Feel free to ignore.
\makeatletter
\def\@makechapterhead#1{
  {\parindent \z@ \raggedright \normalfont
   \Huge\bfseries \thechapter. #1
   \par\nobreak
   \vskip 20\p@
}}
\def\@makeschapterhead#1{
  {\parindent \z@ \raggedright \normalfont
   \Huge\bfseries #1
   \par\nobreak
   \vskip 20\p@
}}
\makeatother

% This macro defines a chapter, which is not numbered, but is included
% in the table of contents.
\def\chapwithtoc#1{
\chapter*{#1}
\addcontentsline{toc}{chapter}{#1}
}

% Draw black "slugs" whenever a line overflows, so that we can spot it easily.
\overfullrule=1mm

%%% Macros for definitions, theorems, claims, examples, ... (requires amsthm package)

\theoremstyle{plain}
\newtheorem{thm}{Theorem}
\newtheorem{lemma}[thm]{Lemma}
\newtheorem{claim}[thm]{Claim}

\theoremstyle{plain}
\newtheorem{defn}{Definition}

\theoremstyle{remark}
\newtheorem*{cor}{Corollary}
\newtheorem*{rem}{Remark}
\newtheorem*{example}{Example}

%%% An environment for proofs

\newenvironment{myproof}{
  \par\medskip\noindent
  \textit{Proof}.
}{
\newline
\rightline{$\qedsymbol$}
}

%%% An environment for typesetting of program code and input/output
%%% of programs. (Requires the fancyvrb package -- fancy verbatim.)

\DefineVerbatimEnvironment{code}{Verbatim}{fontsize=\small, frame=single}

%%% The field of all real and natural numbers
\newcommand{\R}{\mathbb{R}}
\newcommand{\N}{\mathbb{N}}

%%% Useful operators for statistics and probability
\DeclareMathOperator{\pr}{\textsf{P}}
\DeclareMathOperator{\E}{\textsf{E}\,}
\DeclareMathOperator{\var}{\textrm{var}}
\DeclareMathOperator{\sd}{\textrm{sd}}

%%% Transposition of a vector/matrix
\newcommand{\T}[1]{#1^\top}

%%% Various math goodies
\newcommand{\goto}{\rightarrow}
\newcommand{\gotop}{\stackrel{P}{\longrightarrow}}
\newcommand{\maon}[1]{o(n^{#1})}
\newcommand{\abs}[1]{\left|{#1}\right|}
\newcommand{\dint}{\int_0^\tau\!\!\int_0^\tau}
\newcommand{\isqr}[1]{\frac{1}{\sqrt{#1}}}

%%% Various table goodies
\newcommand{\pulrad}[1]{\raisebox{1.5ex}[0pt]{#1}}
\newcommand{\mc}[1]{\multicolumn{1}{c}{#1}}


% Title page and various mandatory informational pages
\begin{document}
%%% Title page of the thesis and other mandatory pages

%%% Title page of the thesis

\pagestyle{empty}
\hypersetup{pageanchor=false}
\begin{center}

\centerline{\mbox{\includegraphics[width=166mm]{img/logo-en.pdf}}}

\vspace{-8mm}
\vfill

{\bf\Large BACHELOR THESIS}

\vfill

{\LARGE{Janika Pokorná}}

\vspace{15mm}

{\LARGE\bfseries{Conjugate Gradient Method for Solving Singular Systems}}

\vfill

Department of Numerical Mathematics

\vfill

{
\centerline{\vbox{\halign{\hbox to 0.45\hsize{\hfil #}&\hskip 0.5em\parbox[t]{0.45\hsize}{\raggedright #}\cr
Supervisor of the bachelor thesis: &RNDr. Jan Papež, Ph.D. \cr
\noalign{\vspace{2mm}}
Study programme: &Mathematical Modeling \cr
\noalign{\vspace{2mm}}
Study branch:&???? \cr
}}}}

\vfill

% Zde doplňte rok
Prague 2024

\end{center}

\newpage

%%% Here should be a bound sheet included -- a signed copy of the "bachelor
%%% thesis assignment". This assignment is NOT a part of the electronic
%%% version of the thesis. DO NOT SCAN.

%%% A page with a solemn declaration to the bachelor thesis

\openright
\hypersetup{pageanchor=true}
\pagestyle{plain}
\pagenumbering{roman}
\vglue 0pt plus 1fill

\noindent
I declare that I carried out this bachelor thesis independently, and only with the cited
sources, literature and other professional sources. It has not been used to obtain another
or the same degree.

\medskip\noindent
I understand that my work relates to the rights and obligations under the Act No.~121/2000 Sb.,
the Copyright Act, as amended, in particular the fact that the Charles
University has the right to conclude a license agreement on the use of this
work as a school work pursuant to Section 60 subsection 1 of the Copyright~Act.

\vspace{10mm}

\hbox{\hbox to 0.5\hsize{%
In \hbox to 6em{\dotfill} date \hbox to 6em{\dotfill}
\hss}\hbox to 0.5\hsize{\dotfill\quad}}
\smallskip
\hbox{\hbox to 0.5\hsize{}\hbox to 0.5\hsize{\hfil Author's signature\hfil}}

\vspace{20mm}
\newpage

%%% Dedication

\openright

\noindent
\Dedication

\newpage

%%% Mandatory information page of the thesis

\openright

\vbox to 0.5\vsize{
\setlength\parindent{0mm}
\setlength\parskip{5mm}

Title:
Conjugate Gradient Method for Solving Singular Systems

Author:
Janika Pokorná

Department:
Department of Numerical Mathematics

Supervisor:
RNDr. Jan Papež, Ph.D.

Abstract:
We rigorously derive the Conjugate Gradient (CG) method for regular systems. Then, using orthogonal decomposition of vectors, we show the conditions necessary for the CG method to work on singular systems and we prove in which cases the CG method will not work at all. We analyze the algorithm mathematically to disprove possible convergence for perturbed right-hand side. We support our results with numerical experiments.

Keywords:
Conjugate Gradients, Singular systems, Orthogonal decomposition

\vss}

\newpage

\openright
\pagestyle{plain}
\pagenumbering{arabic}
\setcounter{page}{1}


%%% A page with automatically generated table of contents of the bachelor thesis

\tableofcontents

%%% Each chapter is kept in a separate file
\chapter*{Introduction}
Finding a solution to a system of linear equations \(\bold{{\mathbb{A}x}=b}\) where the data matrix \(\bold{\mathbb{A}}\) is sparse and has large dimension is a ubiquitous problem in applied mathematics. Many numerical methods have been developed for this purpose, such as \(\bold{LU\,Factorization}\) (which sometimes necessitates pivoting and other algebraic modifications needed to minimize nonzero elements, also called fill-ins) [?], or dimension-reducing methods, such as doplnit. In this thesis we shall focus on the \(\bold{Conjugate\,Gradient\,(CG)}\) method which utilises Krylov subspaces to assure quick convergence and is memory-friendly if applied on systems with certain qualities.


The CG method was formulated in the middle of the 20th century by M. Hestenes and E. Stiefel [\cite{HestenesStiefel52}] as a method for solving regular systems of n linear equations with a sparse symmetrical (or Hermitian, if considering \(A\in{\mathbb{C}^{nxn}}\)) data matrix. Contrary to the often presented popular belief that the CG method is a finite method, this method has been understood (as early as Hestenes and Stiefel) to be iterative, approximating the true result while displaying rounding errors due to the finite precision of computer arithmetic. 
\addcontentsline{toc}{chapter}{Introduction}


\chapter{Conjugate Gradient Method for Regular Systems}

\section{Motivation}

There are numerous ways of approaching the derivation of the Conjugate Gradient method, such as using the Lanczos method (see \cite{vanderVorst03}) or minimizing the quadratic functional \(\mathcal{F}(x)=\frac{1}{2}x^{T}\mathbb{A}x - x^{T}b\). We will be following the second approach, as it offers a basic geometric understanding when derived for \(\mathbb{A}\in{\mathbb{R}^{nxn}}\) and \(b\in\mathbb{R}^n\), \(n\in{N}\). 

\indent To minimize a functional we wish to find its gradient and compare it to zero. Using the definition of gradient, let us compute 

\[\nabla{\mathbb{F}(x)} = \nabla{(\frac{1}{2}x^{T}\mathbb{A}x - x^{T}b)}\]

Linearity of the gradient and a look at the i-th component of the expression gives us the following analysis

\[\nabla{(\frac{1}{2}x^{T}\mathbb{A}x - x^{T}b)}=\frac{1}{2}[\nabla{(x^T\mathbb{A}x)}]_{i} - [\nabla{(x^Tb)}]_i = \frac{1}{2}(\partial_i(x_k\mathbb{A}_{kj}x_j)) - \partial_i(x_kb_k) =\]

\[= \frac{1}{2}(\delta_{ik}\mathbb{A}_{kj}x_j + \delta_{ij}x_k\mathbb{A}_{kj}) -\delta_{ik}x_k = \frac{1}{2}(\mathbb{A}_{ij}x_j + x_k\mathbb{A}_{ki}) - b_i \]


Which gives us 
\begin{equation}\label{1.1}
    \nabla\mathbb{F}(x)=2\frac{1}{2}\mathbb{A}x - b = \mathbb{A}x - b = 0.
\end{equation}

Minimizing the functional \(\mathbb{F}(x)\) thus gives us the solution to the system of equations \(\mathbb{A}x = b\). If we then assume \(\mathbb{A}\) positive definite, such a solution \(x_*\) must be unique.


\section{Derivation of the Conjugate Gradient Method for Regular Systems}


Let us have a sparse positive-definite symmetric matrix \(\mathbb{A}\in\mathbb{R}^{nxn} \) and a right-hand side vector \(b\in\mathbb{R}^n, n\in\mathbb{N} \) large. For the purpose of having a better convergence rate than steepest descent, we shall be minimizing our functional using the \(\bold{"Energy\,Norm"}\), also called A-norm, which we define as 

\[\|x\|_A = \sqrt{(x,x)_A}\] for \(\forall{x}\in\mathbb{R}^n \), where we define the scalar product \[(x,y)_A = x^T\mathbb{A}y \] for \(x,y\in\mathbb{R}^n  \).

This is indeed a norm if \(\mathbb{A} \) is positive definite. Then it satisfies the three necessary conditions:
\begin{itemize}
    \item linearity
    \[\|\alpha x\|_A = \sqrt{\alpha x^T\mathbb{A}\alpha x } = \sqrt{\alpha^2}\sqrt{x^T\mathbb{A} x} = \alpha\|x\|_A, \]
    \item triangle inequality
    \[ (\|x\|_A+\|y\|_A)^2 = {\|x\|_A}^2 + 2\|x\|_A\|y\|_A +  {\|y\|_A}^2\geq {\|x\|_A}^2 + 2x^T\mathbb{A}y +  {\|y\|_A}^2 \]
    \[ = x^T\mathbb{A}x + 2x^T\mathbb{A}y + y^T\mathbb{A}y = (x+y)^T\mathbb{A}(x+y) = {\|x+y\|_A}^2 \forall x,y \in\mathbb{R}^n\]
    using the Cauchy-Schwartz inequality, giving us \(\|x\|_A+\|y\|_A \geq \|x+y\|_A\),
    \item positive-definiteness - stems trivially from positive-definiteness of matrix \(\mathbb{A}\).
\end{itemize}

Let us assume we have an approximation \(x_k\) of the solution \(x_*\), and substitute it into the functional \(\mathbb{F}(x)\):
\begin{equation}
\begin{split}
\mathbb{F}(x_k) = \frac{1}{2}{x_k}^{T}\mathbb{A}x_k - {x_k}^{T}b = \frac{1}{2}{x_k}^{T}\mathbb{A}x_k - {x_k}^{T}\mathbb{A}x_* (\pm\frac{1}{2}{x_*}^{T}\mathbb{A}x_*) \\ = \frac{1}{2}{(x_*-x_k)}^{T}\mathbb{A}(x_*-x_k) - \frac{1}{2}{x_*}^{T}\mathbb{A}x_* = \frac{1}{2}\|x_* - x_k\|_A-\frac{1}{2}\|x_*\|_A
\end{split}
\end{equation}
Since \(x_*\) is a constant value, to minimize the expression on the right we must find \(x_k\) such that \(\|x_* - x_k\|_A\) is minimal. Thus the minimization of the functional \(\mathbb{F}(x_k)\) is equivalent to minimizing the A-norm of \(x_* - x_k \).

Our goal of having a method with quick convergence leads us in the direction of numerical methods on Krylov subspaces (see \cite{LiesenStrakos12}). We would like to find each \(x_k\) from the k-th Krylov subspace \(\mathbb{K}_k(\mathbb{A},r_0), \)
\begin{equation}
r_0 =\mathbb{A}x_0 - b
\end{equation}

If we managed such an iterative method, we would know that the best approximation was found in less that or exactly n steps.

Thus, our method would be an analogy of finding the least squares solution  
\begin{equation}
    \operatorname{arg\,min}_{x_k \in x_0 + \mathbb{K}_k(A,r_0)} \|x_* - x_k\|_A,
\end{equation}
 for k = 1,2,...


- eigenvalues, spectral decomposition and ties to eigenvalues
- notation of solution, estimate, etc
\chapter{CG for Singular Systems}

\section{Solution for right hand side in Span(A)}

\section{Solution for right hand side in Ker(A)}

\section{Analysing the CG algorithm for perturbed right hand side}


\chapter*{Conclusion}
\addcontentsline{toc}{chapter}{Conclusion}


%%% Bibliography
%%% Bibliography (literature used as a source)
%%%
%%% We employ bibTeX to construct the bibliography. It processes
%%% citations in the text (e.g., the \cite{...} macro) and looks up
%%% relevant entries in the bibliography.bib file.
%%%
%%% The \bibliographystyle command selects, which style will be used
%%% for references from the text. The argument in curly brackets is
%%% the name of the corresponding style file (*.bst). Both styles
%%% mentioned in this template are included in LaTeX distributions.

\bibliographystyle{plainnat}    Saad (2003)
\bibliographystyle{unsrt}     [1]

\renewcommand{\bibname}{Bibliography}

%%% Generate the bibliography. Beware that if you cited no works,
%%% the empty list will be omitted completely.

\bibliography{bibliography}

%%% If case you prefer to write the bibliography manually (without bibTeX),
%%% you can use the following. Please follow the ISO 690 standard and
%%% citation conventions of your field of research.

% \begin{thebibliography}{99}
%
% \bibitem{lamport94}
%   {\sc Lamport,} Leslie.
%   \emph{\LaTeX: A Document Preparation System}.
%   2nd edition.
%   Massachusetts: Addison Wesley, 1994.
%   ISBN 0-201-52983-1.
%
% \end{thebibliography}


%%% Figures used in the thesis (consider if this is needed)
\listoffigures

%%% Tables used in the thesis (consider if this is needed)
%%% In mathematical theses, it could be better to move the list of tables to the beginning of the thesis.
\listoftables

%%% Abbreviations used in the thesis, if any, including their explanation
%%% In mathematical theses, it could be better to move the list of abbreviations to the beginning of the thesis.
\chapwithtoc{List of Abbreviations}

%%% Attachments to the bachelor thesis, if any. Each attachment must be
%%% referred to at least once from the text of the thesis. Attachments
%%% are numbered.
%%%
%%% The printed version should preferably contain attachments, which can be
%%% read (additional tables and charts, supplementary text, examples of
%%% program output, etc.). The electronic version is more suited for attachments
%%% which will likely be used in an electronic form rather than read (program
%%% source code, data files, interactive charts, etc.). Electronic attachments
%%% should be uploaded to SIS and optionally also included in the thesis on a~CD/DVD.
%%% Allowed file formats are specified in provision of the rector no. 72/2017.
\appendix
\chapter{Attachments}

\section{First Attachment}

\openright
\end{document}
